\chapter{Folliculitis}

\section*{Definition}
Folliculitis is an inflammation of the hair follicle ostium, typically infectious in origin, presenting as small erythematous papules or pustules with a central hair. It can be superficial (affecting the follicular infundibulum) or deep (involving the entire follicle), and may progress to furuncles or carbuncles.

\section*{Pathophysiology}
Microorganisms colonize and infect the upper hair follicle, triggering neutrophil-mediated inflammation with pustule formation. Occlusion, friction, and moisture disrupt the follicular barrier, facilitating bacterial proliferation. Deep folliculitis extends into the surrounding dermis, producing painful nodular lesions. \textit{Pseudomonas} folliculitis involves endotoxin-mediated inflammation following exposure to contaminated water.

\section*{Etiology}
\begin{itemize}
  \item \textbf{Bacterial:} \textit{Staphylococcus aureus} (most common), \textit{Pseudomonas aeruginosa} (hot tub folliculitis), MRSA
  \item \textbf{Fungal:} \textit{Malassezia} species (Pityrosporum folliculitis), dermatophytes
  \item \textbf{Viral:} Herpes simplex, varicella-zoster (in immunosuppressed)
  \item \textbf{Non-infectious:} Occlusion from tight clothing, shaving, waxing, friction, heavy oils, topical corticosteroids
  \item \textbf{Eosinophilic:} HIV-associated eosinophilic folliculitis, immunosuppressed states
\end{itemize}

\section*{Indications for Evaluation}
Seek care if:
\begin{itemize}
  \item Large, painful nodules or abscess formation (furuncle/carbuncle)
  \item Fever, chills, or rapidly spreading erythema
  \item Recurrent or widespread folliculitis unresponsive to hygiene measures
  \item Immune compromise: diabetes, HIV, chemotherapy, chronic corticosteroid use
\end{itemize}

\section*{Related Symptoms}
\begin{itemize}
  \item Furuncles, carbuncles, impetigo, pseudofolliculitis barbae, hidradenitis suppurativa
\end{itemize}
\textbf{Note:} Hot tub folliculitis caused by \textit{Pseudomonas aeruginosa} typically presents 1--3 days after exposure and resolves spontaneously within 7--10 days without antibiotic treatment.

\section*{Self-Care / Home Management}
\begin{itemize}
  \item Discontinue shaving, waxing, or tight clothing over affected areas until resolved
  \item Cleanse affected areas with benzoyl peroxide 5\% or chlorhexidine wash daily
  \item Warm compresses TID for 15 minutes to promote drainage of pustules
\end{itemize}

\section*{Diagnostic Approach}
\begin{itemize}
  \item Clinical examination noting distribution (beard, scalp, trunk, extremities), morphology, and central hair in pustules
  \item Gram stain and bacterial culture of pustule contents; KOH preparation if fungal etiology suspected
  \item Biopsy with cultures for deep folliculitis or treatment-resistant cases
\end{itemize}

\section*{Treatment / Management Overview}
\begin{itemize}
  \item Mild bacterial: Topical clindamycin 1\% solution/gel BID or mupirocin TID for 5--7 days
  \item Severe or MRSA: Doxycycline 100 mg BID, trimethoprim--sulfamethoxazole DS BID, or cephalexin 500 mg QID
  \item \textit{Malassezia} folliculitis: Topical ketoconazole 2\% shampoo or cream, oral itraconazole for refractory
  \item Incision and drainage for large furuncles or carbuncles; culture-directed antibiotic therapy
\end{itemize}

\clearpage
